\chapter*{2024年\textsf{\textsc{Hdu}}「普通物理2」期中试题}

\vspace{-12.0pt}\setcounter{page}{1}
\begin{center}
  \begin{minipage}{.66\textwidth}
    \centering
    \begin{tasks}
      [ label-width = 5\ccwd, item-indent = 2cm,
        after-item-skip = 0pt ](2)
      \task[\textbf{课程信息}：] 理学院/A0715022
      \task[\textbf{授课教师}：] 陈许敏
      \task[\textbf{解析提供}：] \textsf{Skyrmion}
      \task[{\makebox[5em][s]{\textbf{\textsf\TeX~排版}：}}] \textsf{Skyrmion}
    \end{tasks}
  \end{minipage}
  \begin{minipage}{.15\textwidth}
    \centering
    \begin{tikzpicture}
      \node [ inner sep = 0 ] (wechat) {
        \fancyqr
          [ size = 4em, image padding = 1,
            image = \raisebox{4em/33}{\textcolor{DarkGreen}{\faIcon{weixin}}},
            r color = DarkGreen, version = 4
          ]{http://weixin.qq.com/r/sRJiemTEGjllrSAZ90ej}
        };
      \node [ inner sep = .5ex, align = center, right, font = \scriptsize ]
        at (wechat.east) {未\\央\\学\\社};
    \end{tikzpicture}
  \end{minipage}
  \begin{minipage}{.15\textwidth}
    \centering
    \begin{tikzpicture}
      \node [ inner sep = 0 ] (qq) {
        \fancyqr
          [ size = 4em, image x padding = 2, image y padding = 1,
            image = \raisebox{4em/33}{\textcolor{MidnightBlue}{\faIcon{qq}}},
            l color = MidnightBlue, version = 4
          ]{https://qm.qq.com/q/PH7CsQWvm2}
        };
      \node [ inner sep = .5ex, align = center, left, font = \scriptsize ]
        at (qq.west) {物\\理\\营};
    \end{tikzpicture}
  \end{minipage}
\end{center}

\section{Part A（每空3分，共30分）}

\begin{problem}
  一个小球与轻弹簧组成的系统做简谐运动，若 $t = 0$ 时，小球处于最大位移处，
  则小球简谐运动的初相位是 \result{$0$}.
\end{problem}
\begin{solution}
  $t = 0$ 时 $x = A,~ v = 0$，画出相图比对得 $\phi_0 = 0$.
\end{solution}

\begin{problem}
  一个质点同时参与三个同方向、同频率的简谐振动：
  $ x_1 = 0.2\cos(5t + \pi/6),~ x_2 = 0.2\cos(5t + \pi/3),~
    x_3 = 0.2\cos(5t + \pi/2)$，该质点合振动的振幅为 \result{$0.546$}，
  合振动的初相位为 \result{$\pi/3$}.
\end{problem}
\begin{solution}
  使用旋转矢量法得初相位为 $\pi/3$，合振幅为 $A = 0.2 + \sqrt3 \times 0.2 = 0.546$.
\end{solution}

\begin{problem}
  一横波沿波线传播时的波函数为 $y = 0.02\cos(5\pi t - 2\pi x + \pi/3)$，
  该波的波长为 \result{$\qty1\m$}，波速大小为 \result{$\qty{2.5}{\m/\s}$}.
\end{problem}
\begin{solution}
  $ y = A\cos\ab[\omega\ab(t - \frac xu) + \phi]
      = 0.02\cos\ab[5\pi\ab(t - \frac{x}{2.5}) + \frac\pi3]
  $，波速 $u = \qty{2.5}{\m/\s}$，波长 $\lambda = 2\pi u/\omega = \qty1\m$.
\end{solution}

\begin{problem}
  在牛顿环平凸透镜和平玻璃间充以折射率 $n = 1.33$ 透明液体，第 $8$ 级亮环的直径由液前的
  $\qty{8.80}\mm$ 变成液后的 \result{$\qty{}\mm$}.
\end{problem}
\begin{solution}
  
\end{solution}

\begin{problem}
  $\lambda = \qty{500}\nm$ 的单色光为光源的迈克尔逊干涉实验中，在屏幕中新观测到亮条纹后，
  移动一个反射镜子 $\qty{250}\um$，在屏幕中心飘过 \result{} 个亮条纹.
\end{problem}
\begin{solution}
  
\end{solution}

\begin{problem}
  火车以 $\qty{180}{\km/\hour}$ 的速度行驶，其汽笛的频率为 $\qty{450}\Hz$，
  一个人站在铁轨旁，设空气中的声速为 $\qty{340}{\m/\s}$.
  当火车靠近观察者时，此人听到的汽笛声的频率是 \result{$\qty{527.59}\Hz$}，
  当火车远离观察者时，此人听到的汽笛声的频率是 \result{$\qty{392.31}\Hz$}.
\end{problem}
\begin{solution}
  观察者 (Observer) 静止，所以 $v_o = 0$；发射源 (Source) 接近，所以 $v_s < 0$.
  由多普勒效应公式得观察者所听频率为
  \[
    f = \frac{u \pm v_o}{u \mp v_s}\nu_0
      = \frac{330}{330 - 50} \times \qty{450}\Hz \approx \qty{530.36}\Hz
  \]
  发射源 (Source) 接近时，$v_s > 0$，
  $f = \frac{330}{330 + 50} \times \qty{450}\Hz \approx \qty{390.79}\Hz$.
\end{solution}

\begin{problem}
  当自然光由空气中以 $\qty{45}\degree$ 射向薄膜而反射时，反射光为完全线偏振光，
  折射光相对入射光偏转了 $\qty{15}\degree$，该薄膜的折射率为 \result{$\sqrt2$}
  （空气的折射率为 $1.00$）.
\end{problem}
\begin{solution}
  由折射定律 $n_1\sin i = n_2\sin \gamma$ 得
  $n = \sin\qty{60}\degree/\sin\qty{45}\degree = \sqrt2$.
\end{solution}

\section{Part B （每题10分，共40分）}

\begin{problem}
  一个质量为 $\qty{20}\g$ 的物体作简谐振动，其振幅为 $\qty8\cm$，周期为 $\qty{2.0}\s$.
  当 $t = 0$ 时，位移为 $\qty4\cm$，且沿着正 $x$ 轴方向运动. 试求
  \begin{tasks}[ label = \arabic*. ](2)
    \task 弹簧的倔强系数 $k$.
    \task 简谐振动的位移方程.
  \end{tasks}
\end{problem}
\begin{solution}\leavevmode
  \begin{enumerate}
    \item 由 $T = 2\pi\sqrt{m/k}$ 得 $k = 4\pi^2m/T^2$
  \end{enumerate}
\end{solution}
